\bookchapter{Capitolo 11: Fotografia e gestione delle immagini}

\chapterintro{A dire il vero, la maggior parte delle persone, fotografando con il telefono, ottiene risultati di gran lunga inferiori alle potenzialità del telefono stesso. Non è che tu non sia capace, è che molte piccole impostazioni e abitudini non le hai notate. Questo capitolo non tratta concetti vacui, solo cose che puoi usare subito.}

\booksection{Anche tu commetti questi errori?}

Prima vediamo i problemi più comuni, quanti ne hai collezionati?

\begin{itemize}
  \item \textbf{Lente sporca}: è davvero il problema più diffuso. Quando estrai il telefono, le impronte finiscono sulla lente, le foto escono offuscate e pensi che il telefono sia guasto. Prendi l'abitudine di pulire la lente con un lembo della maglia prima di scattare.
  \item \textbf{Non regoli mai l'esposizione}: in modalità predefinita, toccando lo schermo si mette a fuoco, ma se dopo il tocco scorri verso l'alto o il basso puoi regolare la luminosità. Le foto di molti sono troppo scure o troppo chiare solo perché non conoscono questa operazione.
  \item \textbf{Flash sempre acceso}: le foto con il flash LED del telefono sono davvero brutte, la luce è dura e fredda, le persone sembrano fantasmi. A meno che non sia proprio buio pesto, tienilo spento. In condizioni di bassa luce la modalità notturna è molto meglio del flash.
  \item \textbf{Zoom con due dita}: il pinch sullo schermo è zoom digitale, in pratica ritaglia l'inquadratura e la ingrandisce, con un calo evidente di qualità. Per fotografare soggetti lontani usa i pulsanti di zoom nell'inquadratura (0.5x/1x/2x/3x). Tra questi, 0.5x (grandangolo) e 1x (principale) sono ottiche native, senza perdita di qualità; 2x/3x sono per lo più "zoom a qualità ottica" - realizzato ritagliando il sensore ad alta risoluzione, qualità prossima all'ottica ma non è una lente indipendente. Oltre 3x si tratta in genere di zoom ibrido o digitale, con un calo di qualità evidente.
\end{itemize}

\booksection{Alcuni trucchetti per aumentare subito le buone foto}

Non richiedono conoscenze teoriche, basta usarli una volta per vedere l'effetto:

\begin{itemize}
  \item \textbf{Regola dei terzi + griglia}: nelle impostazioni della fotocamera attiva la "Griglia" (o Linee di riferimento). Inquadrando, posiziona il soggetto da evidenziare in corrispondenza dei punti di intersezione delle quattro linee, e l'immagine risulterà subito "giusta". È la regola più semplice ed efficace della fotografia, niente di esoterico, semplicemente in linea con le abitudini visive dell'occhio umano.
  \item \textbf{Cerca "linee guida"}: quando fotografi paesaggi o scene di strada, fai attenzione alle linee naturali presenti nell'inquadratura - strade, ringhiere, muri, file di alberi - falle arrivare dal bordo dell'inquadratura al soggetto. Usato bene, chi guarda le tue foto penserà che tu abbia studiato fotografia.
  \item \textbf{Non scattare sempre in piedi}: gran parte delle persone scatta sempre all'altezza dei propri occhi, per cui le prospettive sono sempre uguali. Prova a inginocchiarti per scattare dal basso, oppure a cercare un punto più alto per scattare verso il basso, lo stesso posto sembrerà completamente diverso. Per fotografare bambini e animali serve in particolare inginocchiarsi, solo alla loro altezza si ottengono foto naturali.
  \item \textbf{La luce è il primo elemento}: a mezzogiorno con il sole alto, i ritratti hanno ombre dure e brutte sul viso. La luce migliore è la "ora dorata" dell'alba e del tramonto, il sole è obliquo e il colore tende al caldo, tutto viene bene. La luce nuvolosa, pur senza brillare, è uniforme e morbida, per i ritratti è più facile ottenere buoni risultati rispetto al sole pieno.
  \item \textbf{Come usare la modalità notturna}: una volta attivata, il telefono scatta automaticamente una raffica e la fonde. In questo momento la cosa più importante è - la mano non deve tremare assolutamente. Appoggiati a un muro quando puoi, o posa il telefono su un tavolo, se non è possibile trattieni il respiro. La modalità notturna richiede di restare fermi per 2-4 secondi, se tremi l'immagine sfuma.
\end{itemize}

\booksection{Modalità Ritratto e modalità Professionale}

\begin{itemize}
  \item \textbf{Modalità Ritratto}: permette di ottenere lo sfondo sfocato con il soggetto nitido (poca profondità di campo). Poiché le lenti del telefono sono fisicamente piccole, la capacità di sfumatura reale è limitata, quindi ci si affida agli algoritmi AI - in breve, l'AI ritaglia la persona e sfoca lo sfondo. Per l'uso, la distanza migliore è 1-2 metri; troppo lontani o troppo vicini l'algoritmo sbaglia. Inoltre, quando fotografi ritratti, evita linee troppo complesse sullo sfondo, altrimenti il ritaglio AI "fallisce" sui bordi dei capelli.
  \item \textbf{Modalità Professionale}: se vuoi il pieno controllo dell'immagine, prova la modalità professionale. Alcuni parametri chiave da conoscere: ISO più basso = immagine più pulita, ma più scura; tempo di scatto più lento = più luce, ma con la mano che trema sfuma; il bilanciamento del bianco gestisce la temperatura colore, se non è regolato bene la foto tende al blu o al giallo. La maggior parte delle persone non ha bisogno di usare sempre la modalità professionale, ma conoscere questi parametri almeno ti fa capire cosa sta facendo la fotocamera.
\end{itemize}

\booksection{Post-produzione: non esagerare, il giusto basta}

Molti hanno un malinteso sul fotoritocco, pensano che ritoccare non sia reale. In realtà la foto scattata dalla fotocamera è già diversa da ciò che vede l'occhio umano; un'adeguata regolazione serve a restituire ciò che hai visto in quel momento, non a falsificare.

\begin{itemize}
  \item \textbf{Snapseed}: strumento gratuito di fotoritocco di Google, funzioni professionali senza pubblicità. Il mio flusso quotidiano è: prima con "Regola immagine" ritocca luminosità e contrasto con "Dettagli" aumenta la nitidezza con "Bilanciamento bianco" correggi le dominanti cromatiche infine "Ritaglia" per l'inquadratura desiderata. La funzione "Locale" di Snapseed è molto utile, permette di regolare solo una zona dell'immagine senza toccare tutta la foto.
  \item \textbf{Xingtu/Meitu}: queste due sono adatte al ritocco dei ritratti. Ma secondo me c'è un principio: il filtro bellezza entro il 30\%, oltre è finto, la pelle sembra di plastica. Il naturale è più bello del perfetto.
  \item \textbf{Lightroom mobile}: se hai pretese cromatiche, questo è il più professionale. Ha regolazione curve, regolazione HSL dei colori separati e altre funzioni avanzate. Per la maggior parte delle persone però è come usare un cannone per uccidere una mosca.
  \item \textbf{Consiglio sul flusso di ritocco}: la mia abitudine è prima scattare bene l'originale, poi a casa ritoccare con calma sul monitor del computer o sull'iPad; lo schermo del telefono è troppo piccolo e i colori non precisissimi. Se devi proprio ritoccare sul telefono, ricorda di portare la luminosità a un livello medio-alto, non ritoccare in modalità risparmio energetico.
\end{itemize}

\booksection{Troppe foto, come fare? Strategie di organizzazione}

Migliaia di foto nella galleria del telefono, per cercarne una specifica si fa fatica - esperienza comune. Alcuni metodi pratici di organizzazione:

\begin{itemize}
  \item \textbf{Usa la ricerca}: oggi la galleria di iPhone e della maggior parte degli Android supporta la ricerca intelligente. Puoi cercare parole chiave come "gatto", "mare", "documento d'identità", "fiori" e il sistema usa l'AI per riconoscere il contenuto e trovare le foto. Puoi anche cercare per data e luogo, come "maggio dell'anno scorso" o "Pechino". Pochi conoscono questa funzione, ma è davvero utile.
  \item \textbf{Pulizia periodica degli scarti}: screenshot, foto duplicate, scatti sfocati sono i veri colpevoli che occupano spazio. La galleria di iOS ha una funzione automatica di "Elementi duplicati" che li unisce; su Android puoi cercare la cartella "screenshots" nel file manager ed eliminare in blocco gli screenshot non necessari. Io spendo cinque minuti al mese per la pulizia e libero diversi GB.
  \item \textbf{Crea album, non affidarti solo all'ordinamento predefinito}: l'ordinamento cronologico di sistema va bene, ma per contenuti che attraversano più periodi come "crescita del bambino", "foto di viaggio", "documenti importanti" dovresti creare manualmente album per categoria. Gli "Album condivisi" di Apple consentono anche di aggiungere foto insieme alla famiglia: uno crea e tutta la famiglia può vedere.
  \item \textbf{Il backup è il requisito minimo}: le foto sono il dato più insostituibile - i file persi si possono ricostruire, le chat perse sono perse, ma le foto dei bambini, i video degli anziani, una volta persi non ci sono più. Devi assolutamente attivare il backup in cloud. iCloud Foto è il più comodo, su Android ciascun marchio ha il suo servizio cloud, oppure puoi usare il backup automatico di Baidu Netdisk o Alibaba Cloud Drive. Il principio è: almeno una copia in cloud, non mettere tutte le uova nello stesso paniere.
\end{itemize}

\booksection{Screenshot e registrazione schermo}

\begin{itemize}
  \item \textbf{Screenshot}: premi contemporaneamente il tasto accensione + il tasto volume su, è il metodo universale per tutti i telefoni. Molti Android supportano anche lo screenshot con tre dita verso il basso o il doppio tocco con la nocca. Lo screenshot lungo (scroll screenshot) su Android è molto maturo - dopo il primo scatto, premi "Scroll", il telefono scorre automaticamente la pagina e unisce il tutto in un'unica immagine lunga, adatta per screenshot di chat e pagine web.
  \item \textbf{Registrazione schermo}: aggiungi "Registrazione schermo" dal Centro di controllo, tocca per iniziare e tocca di nuovo per fermare. Prima di registrare, premi a lungo il pulsante di registrazione per scegliere se registrare anche l'audio del microfono - se fai un tutorial con commento, attiva il microfono, se registri solo operazioni sullo schermo, disattivalo. Il video registrato viene salvato automaticamente nella galleria.
\end{itemize}

\begin{tcbpractice}{Allenati: completa sul telefono un ciclo completo di scatto + ritocco + organizzazione}
\begin{enumerate}
  \item Apri la "Griglia" nelle impostazioni della fotocamera, trova un luogo con luce naturale e scatta tre foto (ricordati di pulire la lente e di regolare l'esposizione).
  \item Scegli la migliore, usa Snapseed per una regolazione di base (luminosità, contrasto, ritaglio).
  \item Crea un album dedicato per questa foto, chiamalo "Le mie opere".
  \item Verifica che il backup in cloud sia attivo, assicurati che la foto sia già salvata in cloud.
\end{enumerate}

\end{tcbpractice}
